\documentclass[10pt,a4paper]{article}
\usepackage[utf8]{inputenc}\usepackage[T1]{fontenc}\usepackage[german]{babel}
\usepackage{graphicx,float,booktabs,amsmath,siunitx,xcolor,tabularx}
\usepackage[margin=2.2cm]{geometry}\usepackage[hidelinks]{hyperref}
\sisetup{output-decimal-marker={,}}
\hypersetup{pdfauthor={Thomas Schmiegelt},pdftitle={Spulengehäuse eines Stellarators -- vom Datensatz zur Spannungsrechnung}}
\newcommand{\bild}[3]{\begin{figure}[H]\centering\includegraphics[width=#2\linewidth]{bilder/#1}\caption{#3}\end{figure}}
\newcommand{\zwei}[4]{\begin{figure}[H]\centering\includegraphics[width=0.49\linewidth]{bilder/#1}\hfill
  \includegraphics[width=0.49\linewidth]{bilder/#2}\caption{#3}\label{#4}\end{figure}}
\newcommand{\kasten}[2]{\par\smallskip\noindent\fcolorbox{gray}{gray!8}{\parbox{0.96\linewidth}{\textbf{#1} #2}}\par\smallskip}
\title{\textbf{Spulengehäuse eines Stellarators}\\[2pt]\large Vom offenen Datensatz zur Spannungsrechnung --
  ein Projektbericht mit Irrwegen}
\author{Thomas Schmiegelt \textbar{} Hobbyprojekt \texttt{stellarator-design}}
\date{25.10.2026}
\begin{document}\maketitle

\kasten{Hobbyprojekt, aus Neugier entstanden.}{Keine Verbindung zu einem Forschungsinstitut oder Unternehmen,
nicht begutachtet. Alle Zahlen sind Abschätzungen mit vereinfachten Modellen und können Fehler enthalten. Das
W7-X-Modell ist \emph{angelehnt an} Wendelstein~7-X: Die Spulenbahnen stammen aus öffentlichen Daten
\cite{simsopt}; Gehäuse, Verbinder, Verguss, Lagerung und Wandstärken sind eigene, vereinfachte Annahmen und
entsprechen nicht der realen Anlage. Rechnungen, Programmcode und Textentwurf entstanden mit Unterstützung von
KI-Agenten; alle Abbildungen sind eigene Rechnungen und Darstellungen.}

\begin{abstract}
Die supraleitenden Spulen eines Stellarators sind dreidimensional gekrümmte Bauteile, auf die im Betrieb große
magnetische Kräfte wirken. Der Bericht beschreibt, wie in vier Tagen eine durchgängige Rechenkette entstand:
Spulenbahn $\to$ Magnetfeld $\to$ Lorentzkraft $\to$ Gehäusegeometrie $\to$ Vergleichsspannung. Er zeigt auch die
Irrwege: ein falsches Ersatzmodell mit plausiblen Zahlen und ein parametrisches CAD-Modell, das gültig, aber
schwer vernetzbar war. Entscheidend war ein zweiter, paralleler Weg, der die Gehäuse ohne Boolesche Operationen
direkt als Netz erzeugt. Für ein W7-X-ähnliches Spulensystem (70 Spulen, \SI{2,5}{T} auf der Achse) ergeben sich
mit \SI{33}{m^3} Stahl Vergleichsspannungen von 99\,\% unter \SI{43}{MPa} und Verformungen unter \SI{1}{mm}
-- unter stark vereinfachter Lagerung.
\end{abstract}

\tableofcontents

%======================================================================
\section{Worum es geht}
\kasten{Der Stellarator in fünf Sätzen.}{Ein Fusionsplasma (über 100 Mio.\ Grad) wird von einem Magnetfeld in
einem ringförmigen Vakuumgefäß gehalten. Damit die Teilchen nicht nach außen driften, müssen die Feldlinien
schraubenförmig um den Ring laufen. Beim Tokamak erzeugt das ein großer Strom im Plasma selbst; beim
Stellarator ausschließlich die äußeren Spulen, die dafür räumlich verwunden sein müssen. Das macht den
Stellarator im Dauerbetrieb einfacher, die Spulen aber zu anspruchsvollen 3D-Bauteilen mit engen Toleranzen.
Referenz ist Wendelstein~7-X (IPP Greifswald) mit 50 nichtplanaren und 20 planaren supraleitenden Spulen
\cite{bosch2013}.}

Für den Maschinenbau sind die Spulen interessant, weil sich hier Konstruktion, Fertigung und Festigkeit
treffen: Die Bahn ist durch die Physik vorgegeben, die Kräfte durch Strom und Feld, und die Struktur muss
dazwischen Platz finden. Ziel des Projekts war eine nachvollziehbare Rechenkette mit freier Software
(ElmerFEM \cite{elmer}, gmsh \cite{gmsh}, FreeCAD \cite{freecad}) für genau diese Frage.

\subsection{Arbeitsweise}
Gearbeitet wurde mit zwei KI-Agenten: \emph{Hermes Agent} \cite{hermes} (lokal betrieben; Recherche, Projektpflege, lange
Rechenläufe) und \emph{Claude Code} \cite{claudecode} (Programmierung, Fehlersuche, Auswertung). Der Mensch gab Ziel und Richtung
vor, prüfte die Ergebnisse an Bildern und brachte an mehreren Sackgassen die entscheidende konstruktive Idee ein.
Jeder Befund wird mit Beobachtung, Messung, Fundstelle und Status protokolliert; das Protokoll liegt mit dem
Programmcode im Repository.

\begin{table}[H]\centering\small
\begin{tabularx}{\linewidth}{@{}lX@{}}\toprule
Datum & Schritt \\\midrule
22.10. & Projekt angelegt, Recherche zu Funktion, Aufbau und Physik, FEM-Software eingerichtet. \\
23.10. & Validierung am geraden Leiter; parametrisches Ersatzmodell und Bericht v0.1 -- \emph{falsch}; Neustart mit
  echten Spulenbahnen. \\
24.10. & Wickelpakete direkt als Netz, Magnetostatik mit 70 Spulen; Auslegung einer 1-GW-Anlage;
  parametrisches FreeCAD-Gehäuse; Keil und Kollisionen; paralleler Rechenweg und Spannungsrechnung. \\
25.10. & 3D-Schnittmodell mit Leiteraufbau, Visualisierung, dieser Bericht. \\\bottomrule
\end{tabularx}\end{table}

%======================================================================
\section{Recherche: Aufbau einer W7-X-Spule}
Der Ausgangspunkt war Literaturarbeit. W7-X hat einen großen Radius von \SI{5,5}{m}, fünf Feldperioden und
\SI{2,5}{T} auf der Plasmaachse \cite{bosch2013}. Jede Spule ist aus einem Kabel-in-Rohr-Leiter (CICC)
gewickelt: 243 NbTi-Stränge von \SI{0,57}{mm} Durchmesser, jeder mit 144 Filamenten in einer Kupfermatrix
(Cu\,:\,Nicht-Cu = 2,6), in einem koextrudierten Aluminiummantel von \SI{16}{mm} $\times$ \SI{16}{mm} mit \SI{37}{\%} Hohlraum
für das Kühlhelium \cite{rummel2005,maix2006}. Die nichtplanaren Spulen haben 108 Windungen.

Für die Struktur wichtig waren zwei weitere Quellen. Bachmann et al.\ legen für den EU-DEMO die Innenschenkel
der Toroidalfeldspulen als Gewölbe aus, das sich gegenseitig abstützt (\emph{wedged}), mit einer zulässigen
Vergleichsspannung von \SI{660}{MPa} \cite{bachmann2024}. Gil et al.\ nennen für die Spulen von W7-X
Fertigungstoleranzen von \SI{2}{mm} \cite{gil2025}; Schmeling et al.\ zeigen, dass Hochtemperatur-Supraleiter
(REBCO-Band) auf Torsion und Biegung über die harte Achse empfindlich reagieren \cite{schmeling2025}.

%======================================================================
\section{Erste Rechnungen und das falsche Ersatzmodell}
\subsection{Validierung am geraden Leiter}
Vor jeder Spulenrechnung wurde ElmerFEM an einem Fall mit exakter Lösung geprüft: dem unendlich langen Leiter,
$B=\mu_0 I/2\pi r$. Die ersten Läufe lieferten überall null, weil das Netzformat falsch eingelesen und die
Stromdichte mit falschen Schlüsselwörtern angegeben war. Mit dem offiziellen Elmer-Testfall als Vorlage stimmte
das Ergebnis: innerhalb des Leiters \SI{0,09}{\%}, außerhalb \SI{2}{\%} Abweichung (Abb.~\ref{fig:valid}).

\subsection{Das Ersatzmodell (Bericht v0.1)}
Weil die echten Spulenbahnen zunächst als nicht verfügbar galten, wurde ein parametrisches Ersatzmodell aus
Ellipsen und Bändern gebaut und der Strom auf \SI{2,5}{T} kalibriert. Der Vergleich mit W7-X sah gut aus:
Amperewindungen je Spule $0{,}90\times$ des realen Werts.
\zwei{v01_spulen_iso.png}{v01_feld_verteilung.png}{Bericht v0.1: Ersatzmodell und Feldverteilung. Plausible
Zahlen -- aber kein Stellarator.}{fig:v01}

Eine einfache topologische Prüfung zeigte den Fehler: 20 der 50 „nichtplanaren“ Spulen umschlossen die
Plasmaachse gar nicht, die „planaren“ Spulen lagen als waagrechte Bänder über und unter dem Plasma. Die gute
Stromübereinstimmung war Zufall. Die echten Bahnen waren sehr wohl verfügbar: als Fourier-Koeffizienten der
sieben W7-X-Basisspulen in der freien Software simsopt \cite{simsopt}.

Später fanden sich zwei weitere Fehler, die Zahlen rückwirkend änderten: ein Vorzeichenfehler in der Formel der
magnetischen Achse (Kalibrierstrom \SI{12,18}{kA} statt richtig \SI{13,49}{kA} je Windung) und eine vom Löser
selbst gewählte, für 25 gespiegelte Spulen falsche Stromrichtung (\SI{1,38}{T} statt \SI{2,50}{T}).
\kasten{Lehre:}{Plausible Zahlen beweisen nichts. Jede Geometrie wird unabhängig geprüft (Umlaufzahl um die
Achse, Abstände, Feld mit einer zweiten Methode), bevor ein Ergebnis weitergegeben wird.}

%======================================================================
\section{Neustart: Netz direkt aus den Bahndaten}
\subsection{Magnetfeld}
Mit den echten Bahnen (7 Basisspulen, über Dreh- und Spiegelsymmetrie 70 Spulen) zeigte das Feld die bekannten
bohnenförmigen Flussflächen (Abb.~\ref{fig:feld}). Kalibriert auf \SI{2,5}{T} im Mittel auf der Achse ergibt
sich ein Strom von \SI{13,5}{kA} je Windung.
\zwei{w7x_feldlinien_3d.png}{w7x_poincare_bewertung.png}{Feldlinien auf Flussflächen (links), Poincaré-Schnitt
und Verdrillung der Feldlinien über dem Radius (rechts).}{fig:feld}

\subsection{Wickelpakete als Volumennetz}
Der ursprüngliche Plan war eine CAD-Modellierung. Robuster war der direkte Weg: an jedem Bahnpunkt ein
Rechteckquerschnitt (Wickelpaket \SI{192}{mm} $\times$ \SI{160}{mm}, Annahme), je zwei Hälften als Regelfläche,
vernetzt mit gmsh/OpenCASCADE. Sweeps und glatte Lofts scheiterten dagegen an Selbstdurchdringungen der
Oberfläche. Volumenfehler unter \SI{0,4}{\%}; kleinster Spalt zwischen zwei Wickelpaketen \SI{68}{mm}, geprüft
über Mittellinienabstände in \SI{1,5}{s} statt über \SI{45}{min} Boolesche Schnittprüfung.

Die Magnetostatik mit allen 70 Spulen (3,3 Mio.\ Tetraeder, \SI{680}{s}, \SI{10}{GB} RAM) traf das Feld der
Linienleiterrechnung (Biot-Savart) auf der Achse im Median auf \SI{0,24}{\%} (Abb.~\ref{fig:elmer}).
\zwei{w7x_wickelpakete.png}{elmer_70_spulen_vergleich.png}{70 Wickelpakete direkt aus den Bahndaten (links);
FEM gegen Biot-Savart entlang einer Spulennormalen (rechts).}{fig:elmer}
\zwei{validierung_draht.png}{w7x_voxelsolver_validierung.png}{Validierungen: FEM am geraden Leiter (links),
eigener FFT-Feldlöser gegen Biot-Savart (rechts).}{fig:valid}

\subsection{Abstecher: Auslegung auf 1\,GW}
Auf dieselbe Kette setzte eine Studie auf, die vom Ziel her rechnet: \SI{1}{GW} elektrische Leistung $\to$
Plasmagröße und Feld $\to$ eigene Plasmaform (DESC \cite{desc}) $\to$ Spulen (simsopt \cite{simsopt}). Mit
Hochtemperatur-Supraleiter ergibt sich ein großer Radius von etwa \SI{10}{m}. Die Kraft je Spule liegt dann bei
rund \SI{1}{GN}, etwa das Tausendfache von W7-X -- ab hier ist die Tragstruktur die eigentliche
Konstruktionsaufgabe.
\zwei{uebersicht_1GW_A6_A45.png}{struktur_A6_HTS.png}{1-GW-Auslegung (links) und erste Stützschale mit
gekoppelter Rechnung Feld $\to$ Kraft $\to$ Spannung (rechts).}{fig:gw}

%======================================================================
\section{Parametrisches CAD-Modell der Gehäuse}
\subsection{Konstruktionsidee: Variable Section Sweep}
Eine durchgehende Stützschale ließ sich für die kompakte Form nicht sauber konstruieren. Die Idee aus der
Praxis mit Creo war ein \emph{Variable Section Sweep}: je Spule eine geschlossene Leitkurve als Innenkontur des
Gehäuses, Skizzen senkrecht zur Kurve, ein Loft hindurch; alle Wände wachsen nach außen. In FreeCAD wurde das voll
parametrisch umgesetzt: Parametertabelle, Skizzenmaße als Ausdrücke, Wandstärke je Station nach der
Linienlast, Symmetrien als Klone und Verbinder zwischen benachbarten Spulen.
\zwei{freecad_gui_w7x.png}{gehaeuse_A45_HTS_verbunden.png}{Parametrisches FreeCAD-Modell: W7-X-ähnliche
Gehäuse mit Verbindern (links), kompakte 1-GW-Form (rechts).}{fig:fc}

\subsection{Was nicht funktionierte}
\begin{enumerate}\itemsep1pt
\item \textbf{Wandstärken geschätzt statt ausgelegt.} Die erste gekoppelte Rechnung der 1-GW-Form lief durch,
  ergab aber 99\,\% der Vergleichsspannungen unter \SI{5130}{MPa} -- weit über dem Zulässigen. Bei linearer
  Elastizität heißt das: etwa drei- bis fünfmal mehr Stahlquerschnitt.
\item \textbf{Gewölbe (Keil) auf der Innenseite.} In Anlehnung an \cite{bachmann2024} sollten sich die Gehäuse
  innen über Keile abstützen. Das CAD-Modell war gültig, die Vernetzung scheiterte jedoch an deckungsgleichen
  Teilflächen, die aus mehreren Booleschen Schnitten am selben Wickelpaket entstanden. Nach sechs Versuchen
  (Toleranzen, Vereinigen in FreeCAD, bewusste Überlappung) wurde der Keil vorerst herausgenommen.
\item \textbf{Kollisionen.} Zwischen den Skizzenstationen beult ein Loft aus, und an den Ecken greifen schräg
  stehende Nachbarn über. Eine schrittweise Aufbauprüfung (Spule für Spule, mit Kollisionszählung) fand rund
  100 Durchdringungen.
\item \textbf{Lösung im CAD: Nachbarschnitt.} Von jedem Gehäuse werden die Hüllen der Nachbarn plus ein Spiel
  abgezogen. Ergebnis: 0 Kollisionen in allen 70 Aufbauschritten. Die Vernetzung des gesamten Verbunds blieb
  trotzdem aufwendig.
\end{enumerate}
\zwei{struktur_A45_HTS_gehaeuse.png}{gehaeuse_A45_HTS_keil_oben.png}{Erste Spannungsrechnung der FreeCAD-Gehäuse
mit zu dünnen Wänden (links); Gehäuse mit Keil auf der Innenseite -- im CAD gültig, nicht vernetzbar
(rechts).}{fig:keil}
\kasten{Lehre:}{Ein gültiger Volumenkörper ist noch kein vernetzbarer. Je mehr Boolesche Operationen auf
dieselben Flächen treffen, desto eher entstehen Splitterflächen, an denen der Vernetzer scheitert.}

%======================================================================
\section{Paralleler Rechenweg ohne CAD-Kern}
Der Ausweg war, beide Wege zu trennen: das CAD-Modell als Konstruktionsstand, die Rechnung direkt aus den
Bahndaten -- wie bei den Wickelpaketen. Jedes Gehäuse besteht aus zwei Rechteckrohren auf der exakten Bahn
(Außenkontur und Wickelpaket als Aussparung), jeweils als eigenes Oberflächennetz, ganz ohne Boolesche
Operationen. Die Wandstärken folgen aus der Linienlast und dem freien Platz zum Nachbarn. Eine Entflechtung
prüft Oberflächenpunkte fremder Gehäuse (auch zwischen den Stationen) gegen den eigenen Querschnitt und nimmt die
Wand mit dem geringsten Verlust zurück: 4980 $\to$ 40 $\to$ 0 Berührpunkte.

\subsection{Modell und Annahmen}
\begin{table}[H]\centering\small
\begin{tabularx}{\linewidth}{@{}lX@{}}\toprule
Werkstoff Gehäuse & Austenitischer Stahl (316LN), $E=\SI{200}{GPa}$, $\nu=0{,}3$; Vergleichswert
  \SI{667}{MPa} (Annahme: $\tfrac23 R_{p0,2}$ bei \SI{4}{K}) \\
Wickelpaket & homogenisiert (Leiter, Kupfer, Stahl, Isolation), $E=\SI{100}{GPa}$, $\nu=0{,}3$ \\
Last & Lorentzkraft $\vec J\times\vec B$ als Volumenlast, im selben Lauf aus der Magnetostatik (gekoppelt) \\
Lagerung & je Gehäuse ein fest eingespannter Bereich ($r=\SI{0,25}{m}$) am äußersten Bahnpunkt; keine
  Verbinder, kein Tragring \\
Rechnung & linear-elastisch, statisch; Tetraeder, ElmerFEM \\
Nicht modelliert & Verguss, Abkühlung, Vorspannung, Kontakt, Ermüdung \\\bottomrule
\end{tabularx}\end{table}

\subsection{Ergebnis W7-X-ähnlich}
\begin{table}[H]\centering\small
\begin{tabular}{@{}ll@{}}\toprule
Netz & 3,02 Mio.\ Tetraeder, Erzeugung \SI{78}{s} \\
Lösung & \SI{55}{min}, rund \SI{14}{GB} RAM \\
Stahl & \SI{33,4}{m^3} $\approx$ \SI{264}{t} \\
Resultierende Kraft je Spule & FEM \SI{2,42}{MN}, Biot-Savart \SI{2,39}{MN} (+1\,\%) \\
Vergleichsspannung Gehäuse & Mittel \SI{11}{MPa}, 99\,\% $<$ \SI{43}{MPa}, Maximum \SI{110}{MPa} \\
Wickelpakete & Maximum \SI{69}{MPa} \\
Verformung & unter \SI{1}{mm} \\\bottomrule
\end{tabular}\end{table}
Die Spannungen liegen weit unter dem Vergleichswert. Das ist wegen der vereinfachten Lagerung
(Einspannbereiche statt Abstützung über Nachbarn und Tragring) und stellenweise nur \SI{3}{mm} dünner Wände
\emph{kein} Festigkeitsnachweis, aber ein durchgängiger, reproduzierbarer Rechenweg, auf dem sich aufbauen lässt.
\zwei{aufbau_schritt35.png}{struktur_w7x_gehaeuse_netz.png}{Schrittweise Aufbauprüfung im CAD-Ast, Schritt 35/70,
0 Kollisionen (links); Vergleichsspannung im Rechenweg ohne CAD-Kern (rechts).}{fig:aufbau}
\zwei{w7x_feldlinien_spannung.png}{w7x_feldlinien_spannung_geschlossen.png}{Gehäusespannung mit
Magnetfeldlinien auf Flussflächen, aufgeschnitten und geschlossen.}{fig:fls}
\bild{freecad_und_elmer.png}{0.95}{Beide Wege nebeneinander: parametrisches CAD-Modell (links) und
Spannungsrechnung aus dem Rechenweg (rechts).}

%======================================================================
\section{Schichtaufbau im 3D-Modell}
Zum Schluss sollte der innere Aufbau sichtbar werden. Von innen nach außen: Kabel (NbTi-Stränge in Kupfer,
Helium), Aluminiummantel, Windungsisolation, Grundisolierung, Verguss (Quarzsand-Epoxid, eigene Annahme),
Stahlgehäuse; $12\times9=108$ Leiter je Spule.
\zwei{w7x_spule_querschnitt.png}{w7x_torus3d_detail.png}{Schichtaufbau einer Spule (schematisch, links) und im
3D-Schnittmodell (rechts).}{fig:quer}
\zwei{w7x_torus3d_schnitt.png}{w7x_torus3d_sektor.png}{3D-Modell aller 70 Spulen, aufgeschnitten und als
Sektor.}{fig:torus}

\subsection{Wo fließt der Strom?}
Im Betrieb fast vollständig im supraleitenden NbTi, ohne Widerstand. Das Kupfer ist Stabilisator: Es übernimmt
den Strom nur kurzzeitig, wenn der Supraleiter örtlich normalleitend wird (Quench). Aluminiummantel und Helium
tragen keinen Strom; der Mantel gibt dem Leiter Form und Festigkeit und schließt das Helium ein.
\zwei{w7x_leiter3d.png}{w7x_leiter3d_blender.png}{Leiterstück: aufgeschnittener Aluminiummantel, verdrillte
Stränge, daneben ein Strang zwölffach vergrößert mit NbTi-Filamenten (türkis, Anzahl symbolisch).}{fig:leiter}
\bild{w7x_torus3d_blender.png}{0.95}{Gesamtes Spulensystem mit Feldlinien (gerendert in Blender \cite{blender}).}

%======================================================================
\section{Fazit für die Praxis}
\begin{itemize}\itemsep1pt
\item \textbf{Prüfen statt glauben.} Der folgenreichste Fehler sah plausibel aus. Jede Größe wird mit einer
  zweiten Methode gegengerechnet: FEM gegen Biot-Savart, Kraft aus der FEM gegen Kraft aus dem Linienmodell.
\item \textbf{CAD und Berechnungsmodell trennen.} Das parametrische CAD-Modell ist für Konstruktion und
  Weitergabe richtig; für die Rechnung war ein direkt erzeugtes, „defeaturetes“ Netz robuster. Beide parallel zu
  führen hat das Projekt vorangebracht.
\item \textbf{Konstruktive Ideen kamen vom Menschen} (Variable Section Sweep, Nachbarschnitt, paralleler Weg);
  die KI-Agenten haben sie schnell umgesetzt, getestet und dokumentiert.
\item \textbf{Quellen selbst prüfen.} Eine von einem KI-Agenten erstellte Literaturliste enthielt unvollständige
  und falsch zugeordnete Angaben. Für diesen Bericht wurden alle Quellen einzeln gegen Volltext oder
  Verlagsdaten geprüft.
\end{itemize}

\subsection*{Offene Punkte}
Verguss zwischen Wickelpaket und Gehäuse; Verbinder und Tragring statt Einspannbereichen; Mindestwandstärke;
Keil und Gehäuse aus einem Stück; Auslegung der Wandstärken für die 1-GW-Form; Dehnungsnachweis für HTS-Leiter;
Toleranzstudie mit $\pm\SI{2}{mm}$.

\kasten{Fortsetzung (25.10.2026).}{Verguss, Verbinder und Wandstärken -- auch für die 1-GW-Form -- sind inzwischen
in einem Ersatzmodell (Surrogat) aus 157 FEM-Rechnungen untersucht; Bericht
\texttt{bericht/surrogat/surrogat\_bericht.pdf}. Kernaussage: Bei der W7-X-ähnlichen Anlage ist die Seitenwand
der größte Hebel; bei 1~GW reichen Einzelgehäuse mit Verbindern grundsätzlich nicht.}

\bild{gyroid_gefaess_1GW_gesamt.jpg}{1.0}{Ausblick aus dem Surrogat-Bericht: 1-GW-Anlage mit Spulenspannung
(Surrogat, \SI{1000}{MW}), Feldlinien und aufgeschnittenem Plasmagefäß, dessen Wand einen Gyroid-Kühlkern aus
PicoGK hat. Darstellung in Blender, keine Auslegung der Kühlung.}

\subsection*{Software und Daten}
ElmerFEM 9.0 \cite{elmer}, gmsh \cite{gmsh}, FreeCAD 1.1 \cite{freecad}, simsopt \cite{simsopt}, DESC
\cite{desc}, PyVista \cite{pyvista}, ParaView \cite{paraview}, Blender 4.0 \cite{blender}. Programmcode, Befundprotokoll und Abbildungen:
\url{https://github.com/ThomasSchmiegelt/stellarator-design} (MIT-Lizenz).

\begin{thebibliography}{99}\small
\bibitem{bosch2013} H.-S. Bosch, R.\,C. Wolf, T. Andreeva et al.: Technical challenges in the construction of
  the steady-state stellarator Wendelstein 7-X. \emph{Nucl.\ Fusion} 53 (2013) 126001.
  doi:10.1088/0029-5515/53/12/126001
\bibitem{rummel2005} Th. Rummel, K. Rieße, H. Ehmler: Manufacture and test of the non-planar coils for
  WENDELSTEIN 7-X. \emph{Fusion Eng.\ Des.} 75--79 (2005) 117--121. doi:10.1016/j.fusengdes.2005.06.118
\bibitem{maix2006} R.\,K. Maix, V. Bagnato, M. Fricke et al.: Design, production and QA test results of the
  NbTi CIC conductors for the W7-X magnet system. \emph{J.\ Phys.: Conf.\ Ser.} 43 (2006) 753--758.
  doi:10.1088/1742-6596/43/1/184
\bibitem{bachmann2024} C. Bachmann, M. Siccinio, A. Ciula et al.: Re-design of EU DEMO with a low aspect ratio.
  \emph{Fusion Eng.\ Des.} 204 (2024) 114518.
\bibitem{gil2025} P.\,F. Gil, V. Brack, T. Schuler, P. Huslage, E.\,V. Stenson: Manufacturing tolerances of
  non-planar coils for an optimized tabletop stellarator. arXiv:2507.22516 (2025).
\bibitem{schmeling2025} D. Schmeling, M. Russo, B.\,T. Gebreamlak et al.: Experimental results from early
  non-planar NI-HTS magnet prototypes for the Columbia Stellarator eXperiment (CSX). \emph{IEEE Trans.\ Plasma
  Sci.} (2026), doi:10.1109/TPS.2026.3681798; arXiv:2510.07042.
\bibitem{simsopt} M. Landreman, B. Medasani, F. Wechsung, A. Giuliani, R. Jorge, C. Zhu: SIMSOPT: A flexible
  framework for stellarator optimization. \emph{J.\ Open Source Softw.} 6(65) (2021) 3525.
  doi:10.21105/joss.03525. W7-X-Spulendaten: \texttt{simsopt/configs/W7-X.dat}, Spulensatz \texttt{w7x\_v001}.
\bibitem{desc} D.\,W. Dudt, E. Kolemen: DESC: A stellarator equilibrium solver. \emph{Phys.\ Plasmas} 27 (2020)
  102513.
\bibitem{gmsh} C. Geuzaine, J.-F. Remacle: Gmsh: A 3-D finite element mesh generator with built-in pre- and
  post-processing facilities. \emph{Int.\ J.\ Numer.\ Meth.\ Eng.} 79 (2009) 1309--1331.
\bibitem{elmer} ElmerFEM, Version 9.0. CSC -- IT Center for Science, Finnland. \url{https://www.elmerfem.org}
\bibitem{freecad} FreeCAD, Version 1.1. \url{https://www.freecad.org}
\bibitem{pyvista} C.\,B. Sullivan, A. Kaszynski: PyVista: 3D plotting and mesh analysis through a streamlined
  interface for the Visualization Toolkit (VTK). \emph{J.\ Open Source Softw.} 4(37) (2019) 1450.
\bibitem{paraview} J. Ahrens, B. Geveci, C. Law: ParaView: An end-user tool for large data visualization. In:
  \emph{Visualization Handbook}, Elsevier (2005).
\bibitem{blender} Blender, Version 4.0. Blender Foundation. \url{https://www.blender.org}
\bibitem{hermes} Nous Research: Hermes Agent. \url{https://github.com/NousResearch/hermes-agent} (MIT-Lizenz).
\bibitem{claudecode} Anthropic: Claude Code. \url{https://www.anthropic.com/claude-code}
\end{thebibliography}
\end{document}
